\section{The leading constant of the method}\label{sec:optimality}

\emph{Correction (2026-11-27).} An earlier draft of this paper stated that the $\cosh$ kernel is optimal
``within the positivity method'' and that no kernel gives a leading constant below $c^*\approx1.2236$ there.
That was wrong as framed. Proposition~\ref{prop:opt} is correct, but only for positive-definite kernels.
Convolution squares (Proposition~\ref{prop:convsq}) use the same positivity step and do better. The sharp
limit of the method is the constant~$1$ (Proposition~\ref{prop:barrier}).

\subsection{Where the leading coefficient comes from}
Let $K$ be a real kernel and suppose $M$ is a positive-definite function with $\abs{\hat K}\le\hat M$ on $\R$.
Under GRH the zeros $\frac12+i\gamma$ have $\gamma$ real, so the twist $K_y=K(\cdot-y)$ satisfies
$\abs{\sum_\gamma\hat K_y(\gamma)}\le\sum_\gamma\hat M(\gamma)$, and the right side is evaluated by the explicit
formula for $M$. This is the positivity step. For $K=f*\tilde f$ we may take $M=K$; this is the case formalized in
Theorem~\ref{thm:sharp}. In general, only two terms of the resulting bound are large. In the numerator the only
term growing with $q$ is $M(0)\log(q/\pi)$, from the gamma factor through Plancherel. In the denominator the only
term growing with $x$ is $\sqrt x\int K(u)e^{u/2}du$. As $q\to\infty$ and $x/\log^2q\to\infty$ with $K$ fixed,
the method therefore gives
\begin{equation}\label{eq:leading}
  \bias\le\bigl(\rho+o(1)\bigr)\frac{\log q}{\sqrt x},\qquad
  \rho=\frac{M(0)}{\int K(u)e^{u/2}\,du}\ \ge\ \frac{\frac1{2\pi}\int\abs{\hat K}}{\int K(u)e^{u/2}\,du}.
\end{equation}
The last expression is also the leading coefficient in \cite[Lemma~6.1]{lls2015}, which bounds the zero sum by
$\sum_\gamma\abs{\hat K(\gamma)}$ and zero counting instead of by positivity. For $K=f*\tilde f$ with $f$ real,
$\rho=\norm f_2^2/(\langle f,e^{t/2}\rangle\langle f,e^{-t/2}\rangle)$, since $K(0)=\norm f^2$ and
$\hat K(z)=\hat f(z)\overline{\hat f(\bar z)}$. If $\operatorname{supp}K\subseteq(-\infty,T]$, the measure is
supported on $n\le X=xe^{T}$, and the support-normalized constant is $c=e^{T/2}\rho$: then
$\bias\le(c+o(1))\log q/\sqrt X$.

\subsection{Positive-definite kernels}
\begin{lemma}[formal; \lean{GRHPseudo.inner_mul_inner_le}]\label{lem:ranktwo}
In a real inner product space, for all $f,\alpha,\beta$,
\[
  \langle f,\alpha\rangle\langle f,\beta\rangle\le\tfrac12\bigl(\norm\alpha\norm\beta+\langle\alpha,\beta\rangle\bigr)\norm f^2 .
\]
If $\norm\alpha\norm\beta>0$, equality holds iff $f$ is a multiple of $\alpha/\norm\alpha+\beta/\norm\beta$.
\end{lemma}
\begin{proof}
Put $u=\alpha/\norm\alpha$ and $v=\beta/\norm\beta$. Then
$4\langle f,u\rangle\langle f,v\rangle=\langle f,u+v\rangle^2-\langle f,u-v\rangle^2\le\norm f^2\norm{u+v}^2=2(1+\langle u,v\rangle)\norm f^2$.
Equality requires $f\perp u-v$ and $f\parallel u+v$; note that $u+v\perp u-v$.
\end{proof}

\begin{proof}[Proof of Proposition~\ref{prop:opt}]
Work in $L^2[-a,a]$ with $\alpha=e^{t/2}$ and $\beta=e^{-t/2}$. Then $\norm\alpha^2=\norm\beta^2=\int_{-a}^ae^t\,dt=2\sinh a$ and
$\langle\alpha,\beta\rangle=2a$. Lemma~\ref{lem:ranktwo} gives $\rho\ge1/(\sinh a+a)$, with equality iff
$f\propto e^{t/2}+e^{-t/2}=2\cosh(t/2)$ on $[-a,a]$. The factorization statement is the theorem of Boas and Kac
and of Krein: a continuous positive-definite function supported in $[-2a,2a]$ is $f*\tilde f$ for some
$f\in L^2[-a,a]$ \cite{boaskac1945,ehmgneitingrichards2004}. Hence the bound covers every positive-definite
kernel, i.e.\ every kernel that is its own majorant.
\end{proof}

A positive-definite kernel supported in $[-2a,2a]$ has $T=2a$, so its support-normalized constant is at least
$e^a/(a+\sinh a)$, whose minimum over $a$ is $c^*=2/a_*^2\approx1.2236$ (Corollary~\ref{cor:asymp}). Positive-definite
kernels are even. Since the main term weights the window by $e^{u/2}$, an even kernel wastes the lower half of it.

\subsection{Convolution squares}
\begin{proposition}[informal]\label{prop:convsq}
Let $f\ge0$ be real, integrable and supported in $[T/2-L/2,\,T/2]$, and let $K=f*f$. Then $K\ge0$, $\operatorname{supp}K\subseteq[T-L,T]$,
and $M=f*\tilde f$ is a positive-definite majorant: $\abs{\hat K}=\abs{\hat f}^2=\hat M$. Hence the positivity step applies to
$K$ unchanged, with
\[
  \rho=\frac{\norm f_2^2}{\langle f,e^{t/2}\rangle^2},\qquad
  c=e^{T/2}\rho\ \ge\ \frac1{1-e^{-L/2}},
\]
with equality iff $f\propto e^{t/2}$ on its support interval.
\end{proposition}
\begin{proof}
$\hat K=\hat f^2$ and $\hat M=\abs{\hat f}^2$. Next, $M(0)=\norm f^2$ and $\int K(u)e^{u/2}du=\hat K(\frac i2)=(\int f(t)e^{t/2}dt)^2$.
By Cauchy--Schwarz, $\langle f,e^{t/2}\rangle^2\le\norm f^2\int_{T/2-L/2}^{T/2}e^t\,dt=\norm f^2e^{T/2}(1-e^{-L/2})$, with equality iff
$f\propto e^{t/2}$ there.
\end{proof}

At window length $L=4a$, the length of a positive-definite kernel supported in $[-2a,2a]$, this improves on the $\cosh$ kernel:
$1/(1-e^{-2a})$ against $e^a/(a+\sinh a)$, i.e.\ $1.157$ against $1.250$ at $a=1$, $1.084$ against $1.224$ at $a=a_*$, and $1.019$
against $1.313$ at $a=2$. As $L\to\infty$ the constant tends to~$1$, while the positive-definite constant grows. In the
proof of Theorem~\ref{thm:sharp} only the two-sided zero-sum estimate changes: the zero sum of $K_y$ is bounded by that
of $M$ rather than by its own. Theorem~\ref{thm:convclosed} and Corollary~\ref{cor:convcert} are the formal, explicit
version for a smoothed $f=e^{t/2}\mathbf 1_{[-A,0]}$; the Cauchy--Schwarz optimality statement itself is not formalized.

\subsection{The barrier}
\begin{proposition}[informal]\label{prop:barrier}
Let $K\in L^1(\R)$ be real and continuous, with $\hat K\in L^1(\R)$ and $\operatorname{supp}K\subseteq(-\infty,0]$. Then
\[
  \int_\R K(u)e^{u/2}\,du=\frac1{2\pi}\int_\R\hat K(t)\,\frac{\frac12-it}{\frac12+it}\,dt,
  \qquad\text{so}\qquad
  \Bigl\lvert\int_\R K(u)e^{u/2}\,du\Bigr\rvert\le\frac1{2\pi}\int_\R\abs{\hat K(t)}\,dt .
\]
Equality holds for $K(u)=\abs u\,e^{u/2}\mathbf 1_{u\le0}$. Consequently every kernel and every majorant in
\eqref{eq:leading} give support-normalized constant $c\ge1$.
\end{proposition}
\begin{proof}
By Fourier inversion, $K(u)=\frac1{2\pi}\int\hat K(t)e^{itu}dt$ for all $u$. Integrating against $e^{u/2}$ over $u\le0$ and
using Fubini (justified by $\int_{-\infty}^0e^{u/2}du<\infty$ and $\hat K\in L^1$) gives
$\int_{-\infty}^0K(u)e^{u/2}du=\frac1{2\pi}\int\hat K(t)(\frac12+it)^{-1}dt$. Since $K$ is continuous and vanishes on
$(0,\infty)$, $0=K(0)=\frac1{2\pi}\int\hat K(t)\,dt$. Subtracting, and using
$(\frac12+it)^{-1}-1=(\frac12-it)/(\frac12+it)$, which has modulus~$1$, proves the identity and the inequality. For
$K=\abs ue^{u/2}\mathbf 1_{u\le0}$ we have $\hat K(t)=(\frac12-it)^{-2}$, so the integrand is $(\frac14+t^2)^{-1}>0$, and both sides
equal~$1$. For the consequence, translate a kernel supported in $(-\infty,T]$ to $(-\infty,0]$; then
$\int K(u)e^{u/2}du\le e^{T/2}\frac1{2\pi}\int\abs{\hat K}\le e^{T/2}M(0)$, i.e.\ $c\ge1$.
\end{proof}

The extremal kernel is itself a convolution square, $e*e$ with $e(t)=e^{t/2}\mathbf 1_{t\le0}$; its majorant is $e*\tilde e=e^{-\abs u/2}$.
Twisted by $y=\log x$ it gives the weights $\Lambda(n)\log(x/n)/\sqrt x$ on $n\le x$, i.e.\ $\Lambda(n)/\sqrt n$ times
$(n/x)^{1/2}\log(x/n)$. This is a classical Riesz-type weight (compare Bach \cite{bach1990} and \cite[Lemma~2.2]{lls2015}); it is the
member $a=\frac12$ of the family $\Lambda(n)(n/x)^a\log(x/n)$ in Bach's method as presented in \cite{mcgowntrevino2021}.
So the limiting constant $1$ is reached by a classical weight, and the windows of Proposition~\ref{prop:convsq} approach it as $L\to\infty$.
The inequality is a sharp point evaluation for functions whose Fourier transform extends to the upper half-plane. In the one-sided
least-non-residue problem, Carneiro, Milinovich, Quesada-Herrera and Ramos \cite{cmqr2024} find the limiting constant $C(\infty)=1$, with
the same extremal function, by a different (dual) argument. We do not claim novelty for Proposition~\ref{prop:barrier}. Its content for us is
that the constant $1$ is the exact limit of the explicit-formula method for prime-supported small-bias measures. This is a limit of
the method, not a lower bound for the true bias.

\begin{proof}[Proof of Corollary~\ref{cor:asymp}]
\emph{Part 1} is Proposition~\ref{prop:barrier}. \emph{Part 2.} A positive-definite kernel supported in $[-2a,2a]$ has $T=2a$, so
\eqref{eq:leading} and Proposition~\ref{prop:opt} give $c\ge e^a/(a+\sinh a)$. Minimizing:
$(\log c)'(a)=1-(1+\cosh a)/(a+\sinh a)$ vanishes iff $a+\sinh a=1+\cosh a$, i.e.\ $a-1=\cosh a-\sinh a=e^{-a}$. There
\[
  c(a)=\frac{e^a}{1+\cosh a}=\frac{2e^a}{(e^{a/2}+e^{-a/2})^2}=\frac2{(1+e^{-a})^2}=\frac2{a^2}.
\]
Since $c(a)\to\infty$ as $a\to0^+$ and $c(a)\to2$ as $a\to\infty$, while $c(a_*)<2$, the critical point $a_*$ is the global
minimum. \emph{Part 3.} Given $\delta>0$, choose $L$ with $1/(1-e^{-L/2})<1+\delta/2$ and a smooth $f\ge0$ supported in
$[-L/2,0]$, close enough to $e^{t/2}\mathbf 1_{[-L/2,0]}$ in $L^2$ that $c<1+\delta$ in Proposition~\ref{prop:convsq}. For this fixed
kernel every other term of the bound is $O_L(1)$ in the numerator, and $O_L(1)+O(\log q\log x/\sqrt x)$ in the denominator, which is
$o(\sqrt x)$ once $x/\log^2q\to\infty$.
\end{proof}

\subsection{General nonnegative kernels: numerics}
Only $K\ge0$ is needed for the weights to form a measure. Over nonnegative $K$ supported in a window of length $L$, the infimum of
the last expression in \eqref{eq:leading} (normalized by $e^{T/2}$) is a linear program after discretization
(\lean{kernel_lp.py}: piecewise-linear $K$, a polygonal model of $\abs{\cdot}$, solved with HiGHS; we report the exact ratio
of each LP solution). It lies strictly between the barrier and the convolution squares:
\begin{center}
\begin{tabular}{lccccc}
\toprule
$L$ & $4$ & $5.11$ & $6$ & $8$ & $12$\\
\midrule
positive-definite, $e^a/(a+\sinh a)$, $a=L/4$ & $1.250$ & $1.224$ & $1.235$ & $1.313$ & $1.543$\\
convolution squares, $1/(1-e^{-L/2})$ & $1.157$ & $1.084$ & $1.052$ & $1.019$ & $1.003$\\
nonnegative, LP (numerical) & $1.103$ & $1.047$ & $1.024$ & $1.006$ & $1.000$\\
\bottomrule
\end{tabular}
\end{center}
These kernels are not convolution squares, so for them the zero sum has to be bounded through zero counting, as in \cite{lls2015},
and not by positivity. We have not made that route explicit.

\begin{remark}
Lemma~\ref{lem:ranktwo} is formalized. Its specialization to $L^2[-a,a]$, the factorization theorem, Propositions~\ref{prop:convsq}
and~\ref{prop:barrier}, and the asymptotic statement are not. In the one-sided problem of \cite{cmqr2024} the sign conditions are
adapted to a single non-residue condition, and the optimum for the least quadratic non-residue is $0.7615\ldots\log^2p$. Our problem
bounds all characters simultaneously, in both directions.
\end{remark}
